% !TeX root = ../../MainBook.tex
\section{离散时间傅里叶变换}\label{sec:DTFT}

在\refsub{signal}中，我们介绍过以下两种特别的序列：

\begin{definition}[正弦序列]
    $x[n]=\sin(\Omega n)$，其中$\Omega$为数字角频率。
\end{definition}
\begin{definition}[复指数序列]
    $E_{\Omega}[n]=\e^{\im\Omega n}$，其中$\Omega$为数字角频率。
\end{definition}

与\textbf{数字角频率}相对，我们将连续信号的角频率称为\textbf{模拟角频率}。很明显，数字角频率表征了序列的某种频率信息。本章要讨论以下三个问题：
\begin{itemize}
    \item 数字角频率$\Omega$与模拟角频率$\omega$之间有何关联？
    \item 序列何时具有周期性？
    \item 对于周期和非周期的序列，能否像连续信号那样进行分解与合成？
\end{itemize}

第一小节中，我们将解决前两个问题，而最后一个问题是\textbf{离散时间傅里叶变换}(Discrete-Time Fourier Transform, DTFT)与\textbf{离散傅里叶变换}(Discrete Fourier Transform)要解决的，我们将在第二、第三节分别介绍它们。

\subsection{数字角频率}\label{sub:digital_freq}

\subsubsection{数字角频率与模拟角频率的关系}

在学习\refsec{Sampling_Interpolation}后，我们可以另一种视角看待序列，即序列是连续信号的理想采样的结果。已知连续时间信号$f(t)$及采样间隔$T$，则采样得到的序列为$x[n]=f(nT)$。特别地，考虑最典型的具有数字角频率$\Omega$的序列
\begin{equation}
    E_{\Omega}[n]=\e^{\im\Omega n},
\end{equation}
它对应的连续时间信号为
\begin{equation}
    f(t)=\e^{\im\Omega t/T}.
\end{equation}

如果认为序列$x[n]=\e^{\im\Omega n}$具有$\Omega$的数字角频率，则对应的连续信号正是最典型的以$\omega=\Omega/T$为模拟角频率的连续复指数函数。因此，数字角频率与模拟角频率之间的关系为：
\begin{equation}
    \Omega=\omega T.
\end{equation}

可以看到，用连续信号的角频率来理解序列的角频率，则模拟角频率与数字角频率的关系取决于采样间隔$T$。采样间隔越大，振荡模式相同的序列对应的连续信号振荡越缓慢，同样的数字角频率对应的模拟角频率就越小。

然而，前面的讨论忽略了
\begin{equation}
    \e^{\im(\Omega+2k\pi)n}=\e^{\im\Omega n},\quad \forall k\in\mathbb{Z},\ n\in\mathbb{Z}.
\end{equation}

不依赖于连续信号，无法确定序列背后对应着什么模拟角频率。我们在\refsub{alias}中曾遇到过类似的现象，在欠采样的情况下，不同的连续信号会对应同一个序列，这就是\textbf{混叠现象}。

由于数字角频率具有$2\pi$的周期，尽管在定义序列时我们允许任意大小的数字角频率存在，但还是应该在分析序列时限定其范围。本书将以$(-\pi,\pi]$为序列的数字角频率范围\cite{signal_systems}，在一些文献中，也经常以$[0,2\pi)$为数字角频率范围。这两种约定方式没有本质区别。这样，以上两种角频率的关系可以理解为：

\begin{proposition}
    以$T$为采样间隔，模拟角频率为$\omega$的连续信号，采样所得序列的数字角频率为
    \begin{equation}
        \Omega=\omega T+2k\pi,\quad k\in\mathbb{Z},\ \Omega\in(-\pi,\pi].
    \end{equation}
\end{proposition}

混叠是由于对信号的高频成分处理不当而产生的，混叠所得信号一定为（相对）低频的信号，因为此时错误地使用了“恒等式”来进行信号的复原：
\begin{equation}
    \mathcal{F} f=\chi_{[-\nu_m,\nu_m]}\cdot\paren{\mathcal{F} f*\shah_{1/T}},
\end{equation}
其中高于我们所认为的最高频率$\nu_m$的频率成分都被截掉。因此，使用$(-\pi,\pi]$作为数字角频率的范围是合理的，它直接对应着使用序列所能还原出的连续信号的模拟角频率范围$(-\pi/T,\pi/T]$。

\begin{example}
    序列$x[n]=\cos(n\pi)$的数字角频率为$\Omega=\pi$，假定采样间隔$T=1$，对应的连续信号为$f(t)=\cos(\pi t)$，其模拟角频率为$\omega=\pi$。

    利用香农采样定理的方式还原函数，所得连续信号为
    \begin{equation}
        f_{\mathrm{re}}(t)=\sum_{n=-\infty}^{\infty}\cos(n\pi)\sinc\paren{\frac{t-nT}{T}}=\sum_{n=-\infty}^{\infty}(-1)^n\sinc(t-n),
    \end{equation}
    利用$\sin\pi(t-n)=(-1)^n\sin\pi t$，可得
    \begin{equation}
        f_{\mathrm{re}}(t)=\frac{\sin\pi t}{\pi}\sum_{n=-\infty}^{\infty}\frac{(-1)^n}{t-n}=\cos(\pi t).
    \end{equation}
    我们用到了：$\displaystyle\sum_{n=-\infty}^{\infty}\frac{(-1)^n}{t-n}$在主值意义下收敛于$\pi\cot(\pi t)$，可以利用留数定理证明之。

    还原出的函数是以$\Omega\cdot T=\pi$为模拟角频率的函数。然而，$\cos((2k+1)\pi t),\ k\in\mathbb{N}^*$以采样间隔$T=1$采样后也会得到同样的序列，但由于它们的模拟角频率太高，如果不改变采样点，就无法通过采样序列还原出来。
\end{example}

\subsubsection{序列的周期性}

\begin{definition}[序列的周期性]
    序列$x[n]$具有周期性，当且仅当存在正整数$N$，使得
    \begin{equation}
        x[n+N]=x[n],\quad \forall n\in\mathbb{Z}.
    \end{equation}
\end{definition}

带入这个定义，不难发现，尽管$E_{\Omega}[n]=\e^{\im\Omega n}$对应的连续信号$f(t)=\e^{\im\Omega t/T}$总是具有周期性，但$E_{\Omega}[n]$本身未必具有周期性：
\begin{equation}
    E_{\Omega}[n+N]=\e^{\im\Omega(n+N)}=E_{\Omega}[n]\cdot \e^{\im\Omega N}.
\end{equation}
要使$\e^{\im\Omega N}=1$，必须保证$\Omega N=2k\pi,\ k\in\mathbb{Z}$，即
\begin{equation}
    \Omega=2\pi\frac{k}{N}.
\end{equation}
这表明，序列$E_{\Omega}[n]$具有周期性的必要条件是$\Omega/\pi$是有理数；另一方面，当$\Omega/\pi$是有理数时，设$\Omega=\pi\displaystyle\frac{p}{q},\ p,q\in\mathbb{Z}$，且$p/q$为既约分数，容易验证$E_{\Omega}[n]$的最小正周期为$2q$。综上我们得到：

\begin{proposition}\label{prop:discrete_period}
    序列$E_{\Omega}[n]$具有周期性的充要条件是$\Omega/\pi\in\mathbb{Q}$。
\end{proposition}

对于正弦序列$x[n]=A\sin(\Omega n+\varphi)$，它总能表示为复指数序列的线性组合，因此具有相同的周期性判则。

本小节前半部分已经指出，数字角频率$\Omega$的范围为$(-\pi,\pi]$或其他长为$2\pi$的区间，虽然$\Omega=\pi\displaystyle\frac{p}{q}$时$E_{\Omega}[n]$都具有周期性，但真正有效的数字角频率数仅有$2q$个，因为
\begin{equation}
    \Omega=\pi\frac{p}{q}=\pi\frac{p\bmod q}{q}+2k\pi,\quad k\in\mathbb{Z}.
\end{equation}
换言之，\textbf{周期为$N$的序列只有$N$种数字角频率}，我们不妨说这$N$种数字角频率为：
\begin{equation}
    \braces{\Omega_k=\frac{2\pi k}{N}\;\middle|\;k=0,1,\dots,N-1}.
\end{equation}
在\refsec{DFT}中我们将看到，周期为$N$的序列总能表示为具有这$N$种数字角频率的复指数序列的线性组合。

\subsection{离散时间傅里叶变换}\label{sub:dtft}

如何对序列进行分解与合成？具体来说，根据傅里叶级数和傅里叶变换的经验，我们希望能将序列分解为一系列不同数字角频率的\textbf{复指数序列}的组合，以此来表征序列各种频率成分的大小。

在处理连续信号时，我们使用傅里叶级数来分析周期信号，类比到周期序列，需要取定基波数字角频率，考察各次谐波分量的大小。然而，序列的数字角频率有界，需要对这个方法进行一些修正，从而得到周期序列的分析方法，这就是\refsec{DFT}将讨论的\textbf{离散傅里叶变换}(Discrete Fourier Transform, DFT)；对于非周期序列，自然需要类比傅里叶变换，需要考察$(-\pi,\pi]$范围内所有数字角频率成分的大小，并将序列表示为这些频率成分与相应复指数序列的加权积分，这就是本节要介绍的\textbf{离散时间傅里叶变换}(Discrete-Time Fourier Transform, DTFT)。

以$E_{\Omega}[n]=\e^{\im\Omega n}$表示数字角频率为$\Omega$的复指数序列。这样，以上分析非周期序列的想法可以表述为：给定$x[n]$，我们希望找到一个系数函数$\mathcal{X}(\Omega)$，它以$2\pi$为周期，使得
\begin{equation}
    x[n]=\frac{1}{2\pi}\int_{-\pi}^{\pi}\mathcal{X}(\Omega)E_{\Omega}[n]\,\dd\Omega
\end{equation}
（我们很快会看到系数$1/2\pi$的作用）。

为了求出系数函数$\mathcal{X}(\Omega)$，将序列化归为我们熟悉的连续信号：
\begin{equation}
    f(t)\cdot\shah=\sum_{n=-\infty}^{\infty}x[n]\delta_{n}=\sum_{n=-\infty}^{\infty}\paren{\frac{1}{2\pi}\int_{-\pi}^{\pi}\mathcal{X}(\Omega)E_{\Omega}[n]\,\dd\Omega}\delta_{n},
\end{equation}
即
\begin{equation}
    f(t)=\frac{1}{2\pi}\int_{-\pi}^{\pi}\mathcal{X}(\Omega)\e^{\im\Omega t}\,\dd\Omega=\mathcal{F}^{-1}\sbr{\mathcal{X}\cdot\chi_{[-\pi,\pi]}}(t).
\end{equation}

对$f(t)\cdot\shah$做傅里叶变换。由卷积定理，有
\begin{align}
    \mathcal{F}\sbr{f(t)\cdot\shah}(\Omega)&=\frac{1}{2\pi}\mathcal{F} f(\Omega)*\mathcal{F}\shah \nonumber \\
    &=\paren{\mathcal{X}(\Omega)\chi_{[-\pi,\pi]}}*\shah_{2\pi} \nonumber \\
    &=\mathcal{X}(\Omega).
\end{align}
第一个等号用到了\refsub{fourier_of_dirac_comb}中给出的$\shah$的角频率形式傅里叶变换：
\begin{equation}
    {\color{angcolor}\mathcal{F}\shah=\frac{2\pi}{T}\shah_{2\pi/T},}
\end{equation}
最后一个等号用到了$\mathcal{X}$是以$2\pi$为周期的。

另一方面，
\begin{equation}
    f(t)\cdot\shah=\sum_{n=-\infty}^{\infty}x[n]\delta_{n},
\end{equation}
由分布的傅里叶变换的连续性，有：
\begin{equation}
    \mathcal{F}\sbr{f(t)\cdot\shah}(\omega)=\mathcal{F}\sbr{\sum_{n=-\infty}^{\infty}x[n]\delta_{n}}(\omega)=\sum_{n=-\infty}^{\infty}x[n]\e^{-\im\omega n}.
\end{equation}

综上得到：
\begin{align}
    \mathcal{X}(\Omega)&=\sum_{n=-\infty}^{\infty}x[n]\e^{-\im\Omega n}, \\
    x[n]&=\frac{1}{2\pi}\int_{-\pi}^{\pi}\mathcal{X}(\Omega)E_{\Omega}[n]\,\dd\Omega.
\end{align}

由于$\mathcal{X}(\Omega)$可以视作变量$z=\e^{\im\Omega}$的函数，我们令
\begin{equation}
    X(\e^{\im\Omega})=\mathcal{X}(\Omega).
\end{equation}
不难发现，DTFT实际上是z变换在单位圆上的取值，即以上函数$X$的记号实际上就是$x[n]$的z变换的记号：
\begin{equation}
    X(\e^{\im\Omega})=\evalat{X(z)}{z=\e^{\im\Omega}}{}.
\end{equation}

由此可以得到离散时间傅里叶变换的定义：

\begin{definition}[离散时间傅里叶变换]\label{def:DTFT}
    序列$x[n]\in l^1$的DTFT定义为
    \begin{equation}
        X(\e^{\im\Omega})=\sum_{n=-\infty}^{\infty}x[n]\e^{-\im\Omega n},
    \end{equation}
    记$x[n]\xleftrightarrow{\mathrm{DTFT}}X(\e^{\im\Omega})$。其中，$\Omega$为\textbf{数字角频率}，$l^1$见\cref{def:lp}。
\end{definition}

DTFT是良定义的，因为
\begin{equation}
    x[n]\in l^1\implies \sum_{n=-\infty}^{\infty}\abs{x[n]\e^{-\im\Omega n}}\leq \sum_{n=-\infty}^{\infty}|x[n]|<\infty.
\end{equation}

以上探究$x[n]$的分解与合成的过程，已经指出DTFT的反演公式，它可以看作是逆z变换在单位圆上的特例。单位圆在收敛域内，是由$x[n]\in l^1$保证的。
\begin{theorem}[离散时间傅里叶反演公式]\label{thm:DTFT_inv}
    设$x[n]\in l^1$的DTFT为$X(\e^{\im\Omega})$，则$X(\e^{\im\Omega})$的离散时间傅里叶逆变换(IDTFT)即为$x[n]$，即：
    \begin{equation}
        x[n]=\frac{1}{2\pi}\int_{-\pi}^{\pi}X(\e^{\im\Omega})\e^{\im\Omega n}\,\dd\Omega.
    \end{equation}
\end{theorem}

\begin{example}
    求单位冲激序列$\delta[n]$的DTFT。

    解：由定义，
    \begin{equation}
        X(\e^{\im\Omega})=\sum_{n=-\infty}^{\infty}\delta[n]\e^{-\im\Omega n}=1.
    \end{equation}
\end{example}

\begin{example}
    求序列$x[n]=a^n u[n]$的DTFT，其中$|a|<1$。

    解：由等比数列求和公式，
    \begin{align}
        X(\e^{\im\Omega})&=\sum_{n=-\infty}^{\infty}x[n]\e^{-\im\Omega n} \nonumber \\
        &=\sum_{n=0}^{\infty}a^n \e^{-\im\Omega n} \nonumber \\
        &=\sum_{n=0}^{\infty}(a \e^{-\im\Omega})^n \nonumber \\
        &=\frac{1}{1-a \e^{-\im\Omega}}.
    \end{align}

    特别地，当$a=1$时，$x[n]=u[n]\notin l^1$，但不妨写出其DTFT：
    \begin{equation}
        X(\e^{\im\Omega})=\frac{1}{1-\e^{-\im\Omega}}.
    \end{equation}
    可以看到这个函数在$\Omega=0$处发散。

    当$a=\e^{\im\Omega_0}$时，$x[n]=\e^{\im\Omega_0 n}u[n]$，其DTFT为
    \begin{equation}
        X(\e^{\im\Omega})=\frac{1}{1-\e^{\im(\Omega_0-\Omega)}}.
    \end{equation}
    \begin{figure}[H]
        \centering
        \includegraphics[width=0.6\textwidth]{exp_dtft}
        \caption{$a$取不同值时的离散时间傅里叶变换示意图}
    \end{figure}
\end{example}

观察DTFT的定义可知：
\begin{equation}
    \evalat{X(\e^{\im\Omega})}{\Omega=0}{}=\sum_{n=-\infty}^{\infty}x[n],\quad \evalat{X(\e^{\im\Omega})}{\Omega=\pm\pi}{}=\sum_{n=-\infty}^{\infty}(-1)^n x[n].
\end{equation}
也就是说，DTFT在$\Omega=0$处的取值为$x[n]$的总和（或者说直流分量），在$\Omega=\pm\pi$处的取值为$x[n]$的交错和，这两个点处的DTFT值必为实数。上例中，$a$越大，离散时间傅里叶变换在$\Omega=0$处的值越大，在$\Omega=\pm\pi$处的值越小，这符合我们的直观认识：
\begin{equation}
    \sum_{n=0}^{\infty}(-a)^n=\frac{1}{1+a}
\end{equation}
关于$a$单调递减。

\subsection{离散时间傅里叶变换的运算性质}\label{sub:property_of_dtft}

DTFT的运算性质列举如下：
\begin{proposition}[离散时间傅里叶变换的运算性质]
    \label{prop:DTFT_operation}
    设$x[n],y[n]\in l^1$，其DTFT分别为$X(\e^{\im\Omega}),Y(\e^{\im\Omega})$，则有：
    \begin{itemize}
        \item 线性：$\forall a,b\in\mathbb{C}$，有$a x[n]+b y[n]\xleftrightarrow{\mathrm{DTFT}} a X(\e^{\im\Omega})+b Y(\e^{\im\Omega})$
        \item 时移性质：$x[n-n_0]\xleftrightarrow{\mathrm{DTFT}}\e^{-\im\Omega n_0}X(\e^{\im\Omega})$
        \item 频移性质：$\e^{\im\Omega_0 n}x[n]\xleftrightarrow{\mathrm{DTFT}}X(\e^{\im(\Omega-\Omega_0)})$
        \item 共轭性质：$\overline{x[n]}\xleftrightarrow{\mathrm{DTFT}}\overline{X(\e^{-\im\Omega})}$
        \item 反转性质：$x^-[n]\xleftrightarrow{\mathrm{DTFT}}X(\e^{-\im\Omega})$
        \item 时域卷积和定理：$x[n]*y[n]\xleftrightarrow{\mathrm{DTFT}}X(\e^{\im\Omega})Y(\e^{\im\Omega})$
        \item 频域卷积定理：$x[n]y[n]\xleftrightarrow{\mathrm{DTFT}}\displaystyle\frac{1}{2\pi}\int_{-\pi}^{\pi}X(\e^{\im\theta})Y(\e^{\im(\Omega-\theta)})\,\dd\theta$
        \item 频域微分性质：$nx[n]\xleftrightarrow{\mathrm{DTFT}}\im\displaystyle\frac{\dd X(\e^{\im\Omega})}{\dd\Omega}$
    \end{itemize}
\end{proposition}
\begin{proof}
    利用
    \begin{equation}
        X(\e^{\im\Omega})=\evalat{X(z)}{z=\e^{\im\Omega}}{},
    \end{equation}
    以上性质除最后一条以外全部是z变换的性质在单位圆上的特例。这里仅给出频域微分性质的证明。由于
    \begin{equation}
        \abs{X(\e^{\im\Omega})}=\abs{\sum_{n=-\infty}^{\infty}x[n]\e^{-\im\Omega n}}\leq \sum_{n=-\infty}^{\infty}|x[n]|=\norm{x[n]}_1<\infty,
    \end{equation}
    $X(\e^{\im\Omega})$作为一个函数项级数是一致收敛的，因此可以逐项求导：
    \begin{equation}
        \im\frac{\dd X(\e^{\im\Omega})}{\dd\Omega}=\im\sum_{n=-\infty}^{\infty}x[n]\frac{\dd}{\dd\Omega}\e^{-\im\Omega n}=\sum_{n=-\infty}^{\infty}nx[n]\e^{-\im\Omega n}.
    \end{equation}
    注意它与$z$域微分性质（\cref{prop:z_operation}）的形式不同。
\end{proof}

作为共轭性质的推论，我们有：
\begin{corollary}
    如果$x[n]$为实数序列，则$X(\e^{\im\Omega})$具有共轭对称性：$X(\e^{\im\Omega})=\overline{X(\e^{-\im\Omega})}$，即$|X(\e^{\im\Omega})|$为偶函数，$\arg X(\e^{\im\Omega})$为奇函数。
\end{corollary}

频域卷积定理的名称可以这样理解：如果记
\begin{equation}
    \mathcal{X}(\Omega)=X(\e^{\im\Omega}),\quad \mathcal{Y}(\Omega)=Y(\e^{\im\Omega}),
\end{equation}
则右式是周期函数的卷积（\cref{def:period_conv}）：
\begin{equation}
    x[n]y[n]\xleftrightarrow{\mathrm{DTFT}}\frac{1}{2\pi}\mathcal{X}(\Omega)*\mathcal{Y}(\Omega)=\frac{1}{2\pi}\int_{-\pi}^{\pi}\mathcal{X}(\theta)\mathcal{Y}(\Omega-\theta)\,\dd\theta.
\end{equation}

此外，我们给出DTFT的\textbf{帕塞瓦尔定理}(Parseval's Theorem)，它表明时域能量与频域能量成正比：
\begin{theorem}[离散时间傅里叶变换的帕塞瓦尔定理]\label{thm:DTFT_plancherel}
    设$x[n],y[n]\in l^1\cap l^2$，其DTFT分别为$X(\e^{\im\Omega}),Y(\e^{\im\Omega})$，则有
    \begin{equation}
        \sum_{n=-\infty}^{\infty}x[n]\overline{y[n]}=\frac{1}{2\pi}\int_{-\pi}^{\pi}X(\e^{\im\Omega})\overline{Y(\e^{\im\Omega})}\,\dd\Omega.
    \end{equation}
    特别地，当$x[n]=y[n]$时，有
    \begin{equation}
        E=\sum_{n=-\infty}^{\infty}|x[n]|^2=\frac{1}{2\pi}\int_{-\pi}^{\pi}|X(\e^{\im\Omega})|^2\,\dd\Omega,
    \end{equation}
    其中$E$为序列$x[n]$的能量。
\end{theorem}
\begin{proof}
    由离散时间傅里叶反演公式，有
    \begin{align}
        \sum_{n=-\infty}^{\infty}x[n]\overline{y[n]}&=\sum_{n=-\infty}^{\infty}x[n]\cdot\frac{1}{2\pi}\int_{-\pi}^{\pi}\overline{Y(\e^{\im\Omega})}\e^{-\im\Omega n}\,\dd\Omega \nonumber \\
        &=\frac{1}{2\pi}\int_{-\pi}^{\pi}\overline{Y(\e^{\im\Omega})}\paren{\sum_{n=-\infty}^{\infty}x[n]\e^{-\im\Omega n}}\,\dd\Omega \nonumber \\
        &=\frac{1}{2\pi}\int_{-\pi}^{\pi}X(\e^{\im\Omega})\overline{Y(\e^{\im\Omega})}\,\dd\Omega.
    \end{align}
    积分与求和的交换合法，是因为绝对可和能够保证以上函数项级数的一致收敛性：
    \begin{equation}
        \forall\Omega\in(-\pi,\pi],\quad \abs{\sum_{n=-\infty}^{\infty}x[n]\overline{Y(\e^{\im\Omega})}\e^{-\im\Omega n}}\leq\abs{Y(\e^{\im\Omega})}\sum_{n=-\infty}^{\infty}|x[n]|<\infty.
    \end{equation}
\end{proof}

利用帕塞瓦尔定理，可以得到原函数与采样序列的能量关系：

\begin{theorem}
    设$f(t)\in L^1(\mathbb{R})\cap L^2(\mathbb{R})$是带限函数，采样频率高于奈奎斯特频率，以$T$表示采样间隔，$E_d$表示采样序列$x[n]$的能量，$E$表示原函数$f(t)$的能量，则它们满足以下关系：
    \begin{equation}
        E=TE_d,\quad \text{即}\ \sum_{n=-\infty}^{\infty}|x[n]|^2=\frac{1}{T}\norm{f(t)}_2^2.
    \end{equation}
\end{theorem}
\begin{proof}
    在采样频率高于奈奎斯特频率时不会发生混叠，即频域上周期化时不产生交叠，因此离散时间傅里叶变换$X(\e^{\im\Omega})$在$(-\pi,\pi]$范围内就是原函数的傅里叶变换$\mathcal{F} f(\omega)$在$(-\pi/T,\pi/T]$范围内的伸缩，即
    \begin{equation}
        X(\e^{\im\Omega})=\mathcal{F} f(\Omega/T),\quad \Omega\in(-\pi,\pi].
    \end{equation}
    因此有
    \begin{align}
        E_d&=\sum_{n=-\infty}^{\infty}|x[n]|^2 \nonumber \\
        &=\frac{1}{2\pi}\int_{-\pi}^{\pi}|X(\e^{\im\Omega})|^2\,\dd\Omega \nonumber \\
        &=\frac{1}{2\pi}\int_{-\pi}^{\pi}|\mathcal{F} f(\Omega/T)|^2\,\dd\Omega \nonumber \\
        &=\frac{1}{T}\cdot\frac{1}{2\pi}\int_{-\pi/T}^{\pi/T}|\mathcal{F} f(\omega)|^2\,\dd\omega \qquad (\omega=T\Omega) \nonumber \\
        &=\frac{1}{T}\norm{f(t)}_2^2=\frac{E}{T}.
    \end{align}
\end{proof}